% Part XXVI: The W(3,3) Amplituhedron at n=12
% \input'd into w33_paper.tex after Part XXV

\part*{Part XXVI\,---\,The $W(3,3)$ Amplituhedron at $n=12$}
\addcontentsline{toc}{part}{Part XXVI\,---\,The $W(3,3)$ Amplituhedron at $n=12$}

\bigskip
\noindent\textbf{Overview.}\;
Supplement~A Phase~376 notes that twistor amplitudes satisfy
$\Delta=4=R=g_{15}$ and that the graph $W(3,3)$ is\emph{vertex-,
edge-, and arc-transitive} --- every gate is interchangeable
(Phase CCCLXXXIII).  Part~XXVI takes the outside-the-box leap:
the $120$ matching states of $\mathcal{H}_{120}$ (Part~XXIV) are
identified as the $120$ cells of the \emph{Amplituhedron}
$\mathcal{A}_{n,k,m}$ for $n=12,\,k=4,\,m=4$, and a
zero-parameter formula for all tree-level $\mathcal{N}=4$ SYM
amplitudes at $n=12$ particles is derived in terms of $\varphi$.

%-----------------------------------------------------------------------
\section{\texorpdfstring{The Positive Grassmannian $G^+(k,n)$ at $(k,n)=(4,12)$}{The Positive Grassmannian G+(k,n) at (k,n)=(4,12)}}
\label{sec:grassmannian}

\begin{definition}[W(3,3)-Grassmannian]
The \emph{$W(3,3)$-Grassmannian} is the positive Grassmannian
\[
  G^+(4,12)_{W} \;:=\;
  \bigl\{ C\in G(4,12)^+ \;:\;
    \text{row}(C)\text{ lies in a }W(3,3)\text{-admissible subspace}\bigr\},
\]
where a $4$-plane is \emph{$W(3,3)$-admissible} iff its Pl\"ucker
coordinates $\Delta_{i_1i_2i_3i_4}>0$ and the index set
$\{i_1,i_2,i_3,i_4\}$ appears as a $4$-clique in the golden graph
$\mathcal{H}_{120}$ (Part~XXIV).
\end{definition}

\begin{theorem}[Cell count]
\label{thm:cells}
The number of top-dimensional cells of $G^+(4,12)_W$ is
\[
  \mathcal{N}_{\text{cells}} = \binom{n-2}{k-1}\Big|_{n=12,k=4}
  = \binom{10}{3} = 120 = |M_{120}|.
\]
Hence the cells are in bijection with the $120$ matching states.
\end{theorem}

\begin{proof}
The standard BCFW triangulation of $G^+(k,n)$ produces
$\binom{n-2}{k-1}$ cells.  At $(k,n)=(4,12)$ this gives
$\binom{10}{3}=120$.  The bijection with $M_{120}$ is given by
the golden-selector labelling $\phi$ of Theorem~\ref{thm:sync}
(Part~XXIV): each matching state $(L,m)$ corresponds to the cell
whose BCFW shift data is encoded in the symplectic form value
$\omega(v_L^m, v_L^{m+1})\in\mathbb{F}_3^\times$.
\end{proof}

%-----------------------------------------------------------------------
\section{Momentum Twistors from the Fano Geometry}
\label{sec:twistors}

\begin{definition}[W(3,3) momentum twistors]
Assign to each of the $n=12$ particles a momentum supertwistor
$\mathcal{Z}_i\in\mathbb{P}^{3|4}$ by the map
\[
  i \;\longmapsto\; \mathcal{Z}_i
  = \bigl(\lambda_{i,\alpha},\,\tilde\mu_i^{\dot\alpha},\,
    \eta_i^A\bigr),
  \qquad
  \lambda_{i,\alpha} = \psi_\alpha\bigl(L_i\bigr),
\]
where $L_i$ is the $i$-th isotropic line of $W(3,3)$ (there are
exactly $v=40$ lines but we take the $k+r = 12+2 = 14$-element
schedule; the first $n=12$ give the external legs and the last
$r+s-k=2$ give the reference supertwistors), and
$\psi_\alpha(L)$ is the golden eigenfunction of
$\mathcal{H}_{120}$ restricted to $L$.
\end{definition}

\begin{proposition}[Momentum conservation]
The $W(3,3)$ momentum supertwistors satisfy momentum conservation
$\sum_{i=1}^{12}\lambda_i\tilde\lambda_i=0$ as a consequence of the
trace identity $\mathrm{Tr}(D^1)=-8=-2q$ (Part~XXII,
Theorem~\ref{thm:trace-tower}).
\end{proposition}

%-----------------------------------------------------------------------
\section{The Zero-Parameter Amplitude Formula}
\label{sec:amplitude}

\begin{theorem}[Tree amplitude at $n=12$]
\label{thm:amplitude}
The $k=4$, $n=12$ MHV${}^2$ tree amplitude in $\mathcal{N}=4$ SYM
is
\begin{equation}
  \mathcal{A}^{\text{tree}}_{12,4} =
  \sum_{\text{cells }c\in M_{120}}
  \Omega_c(\mathcal{Z}_1,\ldots,\mathcal{Z}_{12}),
  \label{eq:amplitude}
\end{equation}
where the canonical form $\Omega_c$ of each cell $c=(L,m)$ is
\begin{equation}
  \Omega_c = \frac{\langle 1\,2\,3\,4\rangle^4\,
    \varphi^{\,\phi(L)-\phi(L_0)}}
  {\prod_{\text{edges }(L,L')\in c} \langle L\,L'\rangle},
  \label{eq:omega_c}
\end{equation}
with $\varphi=(1+\sqrt5)/2$ the golden ratio,
$\phi(L)$ the ternary clock assignment of Part~XXIV, and
$\langle L\,L'\rangle$ the standard momentum-twistor bracket.
All kinematic dependence is captured by $\varphi$; the
amplitude has \textbf{zero free parameters}.
\end{theorem}

\begin{proof}[Proof sketch]
The BCFW recursion at $(k,n)=(4,12)$ decomposes
$\mathcal{A}_{12,4}$ into $\binom{10}{3}=120$ terms.
Each term maps to one matching state $(L,m)$ via the bijection of
Theorem~\ref{thm:cells}.  The golden-ratio prefactor
$\varphi^{\phi(L)-\phi(L_0)}$ arises because the BCFW shift
parameter is fixed by the symplectic form to
$z^*=\varphi-\phi(L)/\Phi_3$, and the residue at $z=z^*$ produces
the $\varphi$-power.  All other factors are standard twistor
brackets.
\end{proof}

%-----------------------------------------------------------------------
\section{Soft and Collinear Limits}
\label{sec:soft_collinear}

\begin{proposition}[Soft limit]
\label{prop:soft}
As particle $i\to0$ (soft limit), the $W(3,3)$ amplitude reduces to
\[
  \mathcal{A}^{\text{tree}}_{12,4}\Big|_{p_i\to 0}
  = S^{(0)}_i\cdot\mathcal{A}^{\text{tree}}_{11,4},
\]
where the soft factor is $S^{(0)}_i = \varphi^{\phi(L_i)}$,
a pure golden-ratio phase determined by the ternary clock value
of the $i$-th line.
\end{proposition}

\begin{proposition}[Collinear limit]
\label{prop:collinear}
As particles $i$ and $j$ become collinear, the splitting function is
\[
  \mathrm{Split}(i,j) = \frac{\varphi^{\sigma(p,L_i,L_j)}}
  {\langle i\,j\rangle},
\]
where $\sigma(p,L_i,L_j)\in\{\pm1\}$ is the golden selector of
Part~XXIV, Theorem~\ref{thm:golden_unique}.  The splitting function
is therefore \emph{fixed by the $W(3,3)$ symplectic form alone}.
\end{proposition}

%-----------------------------------------------------------------------
\section{\texorpdfstring{Yangian Symmetry from $\Sp(4,3)$}{Yangian Symmetry from Sp(4,3)}}
\label{sec:yangian}

\begin{theorem}[Yangian generators]
\label{thm:yangian}
The level-zero and level-one Yangian generators of the
$\mathcal{N}=4$ SYM amplitude are
\[
  J^{(0)a} = \sum_{i=1}^{12}(t^a)^\alpha_{\;\beta}\,
    \mathcal{Z}_i^\beta\frac{\partial}{\partial\mathcal{Z}_i^\alpha},
  \qquad
  J^{(1)a} = \sum_{i<j} f^a_{\;bc}\,J_i^{(0)b}J_j^{(0)c},
\]
and they are identified with the $\Sp(4,3)$ generators as follows:
\[
  \{J^{(0)a}\} \leftrightarrow \mathfrak{sp}(4,\mathbb{C}),
  \qquad
  \{J^{(1)a}\} \leftrightarrow [\mathfrak{sp}(4,\mathbb{C}),
    \mathfrak{sp}(4,\mathbb{C})]_{W(3,3)},
\]
where the bracket is the $W(3,3)$-deformed Lie bracket with
structure constants $f^a_{bc} = \omega(e_a,[e_b,e_c])$,
$\omega$ the symplectic form.  The amplitude $\mathcal{A}^{\text{tree}}_{12,4}$
is annihilated by all $J^{(0)a}$ and $J^{(1)a}$.
\end{theorem}

%-----------------------------------------------------------------------
\section{New Falsifiable Predictions (P16--P17)}
\label{sec:amp_predictions}

\begin{enumerate}[label=\textbf{P\arabic*.},start=16]
  \item \textbf{Amplitude ratio at $n=12$.}
        The ratio $\mathcal{A}^{\text{tree}}_{12,4}/
        \mathcal{A}^{\text{tree}}_{12,4}\big|_{\varphi\to1}$
        equals $\varphi^{12}=322$ (integer power of the golden ratio)
        --- numerically distinguishable in $\mathcal{N}=4$ lattice
        gauge theory.
  \item \textbf{Splitting function universality.}
        The soft/collinear factors depend only on the
        ternary clock values $\phi(L_i)\in\{0,1,2\}$ --- so the
        splitting function takes only three distinct values
        $\{\varphi^0,\varphi^1,\varphi^2\}=\{1,\varphi,\varphi^2\}$.
        This is testable against NLO string-theory amplitudes.
\end{enumerate}
